\section{Lessons Learned}

\textbf{Observability is a practice, not an installation.} The dashboards, rules, and labels that
make sense today will not all make sense in six months. The system changes, the questions change,
and the platform has to follow. Treating it as something you set up once and leave alone is how it
stops being useful.

\textbf{GitOps gives the platform memory} \cite{lsst-it-k8s-cookbook}. When a dashboard changes,
there is a reason, a review, and a record. When an alert is added or removed, the decision is
visible to anyone who looks. That history turns out to be more valuable than it seems at the time,
especially after an incident, when understanding what changed and why is exactly the question you
need to answer.

\textbf{Labels are design decisions, not metadata}
\cite{prometheus-operator,grafana-loki}. Good labels make it possible to correlate signals across
systems. Bad labels make queries slow, expensive, or just confusing. Label discipline is an
engineering interface between teams, and it deserves the same kind of care as an API.

\textbf{Seeking help from a contractor is always a good practice} when there is uncertainty about
how to move forward; however, the contractor needs very specific instructions on what to do. That
is not simple when there is little knowledge of the tools' capacity. It feels like another
contractor is needed to define the requirements to be delivered to the observability contractor.
The most reasonable approach is to give some liberty to the contractor to recommend what is best
based on their experience while maintaining very tight control over the project's progress. There
will be many situations in which requirements will be changed, added, or removed, and the
contractor needs to be aware of this. One decision that made a real difference was asking the
contractor to work alongside us rather than just deliver something. The co-work sessions gave the
team a chance to absorb how the platform was being built, why certain decisions were made, and
where the tricky parts were. By the time the contract ended, we were not starting from scratch; we
understood what we had inherited and were ready to take it further. That knowledge transfer was
probably worth as much as the platform itself.

\textbf{Dashboard work is human work.} The best dashboard is not the one with the most panels. It
is the one that helps the right person answer a question when they are tired, under pressure, and
do not have time to think. That is a harder design problem than it looks.

\textbf{Open source is a good fit when the team is ready to own it.} The transparency and
flexibility are real advantages. So is the cost: someone has to understand the upgrades, the
sizing, the edge cases, and the failure modes. For an observatory of this magnitude, that ownership is not a burden; it is part of what makes the platform
sustainable.